% !TEX root = ../00_Main/Main.tex
\documentclass[../00_Main/Main.tex]{subfiles}
\begin{document}

\label{sec:results}

Los valores a los que se hace referencia, salvo que se especifique otra
cosa, corresponden al desempe\~no normalizado $\bar r$
(Ecuaci\'on~\eqref{eq:normalised-severity-materials}) medio de cada
escenario; \emph{referencia} es la condici\'on de entrenamiento y
\emph{generalizaci\'on}, la media en los $21$ escenarios. Los mapas
estad\'isticos por escenario y las matrices de comparaci\'on entre pares
pueden encontrarse en el Ap\'endice~\ref{ap:stat-maps}
(Figuras~\ref{fig:heatmap-welch}--\ref{fig:ensemble-pairwise}).

\subsection{Mejor modelo individual}\label{sec:res-individual}

El Transformer demuestra ser el mejor modelo individual
(Tabla~\ref{tab:global-mean}): $0{,}927$ en la referencia y $0{,}930$ en
generalizaci\'on, frente a $0{,}871$--$0{,}899$ del resto de los modelos
optimizados, con la menor dispersi\'on entre secuencias y la mayor media
en $16$ de los $22$ escenarios (Figura~\ref{fig:heatmap-global}). Su
ventaja es de tama\~no medio frente a DQN ($d = 0{,}51$) y al DQN
recurrente ($d = 0{,}57$), y medio frente a Q-Learning ($d = 0{,}78$).

\begin{table}[H]
\centering
\caption[Desempe\~no de los modelos individuales]{Modelos individuales:
media en la referencia, desviaci\'on est\'andar entre secuencias
($\sigma$), media en generalizaci\'on, peor escenario y p\'erdida
respecto de la referencia.}
\label{tab:global-mean}
\small
\setlength{\tabcolsep}{4pt}
\begin{tabular}{@{}lccclc@{}}
\toprule
\textbf{Modelo} & \textbf{Ref.} & $\sigma$ & \textbf{Gen.} & \textbf{Peor escenario} & \textbf{P\'erdida} \\
\midrule
Transformer       & $0{,}927$ & $0{,}045$ & $0{,}930$ & Secuencias largas ($0{,}860$) & $0{,}068$ \\
DQN               & $0{,}894$ & $0{,}055$ & $0{,}899$ & Secuencias largas ($0{,}841$) & $0{,}053$ \\
A2C               & $0{,}887$ & $0{,}063$ & $0{,}896$ & Secuencias largas ($0{,}826$) & $0{,}062$ \\
DQN recurrente    & $0{,}899$ & $0{,}049$ & $0{,}889$ & Extrap.\ de severidad ($0{,}831$) & $0{,}068$ \\
Double Q-Learning & $0{,}896$ & $0{,}048$ & $0{,}877$ & Extrap.\ de severidad ($0{,}785$) & $0{,}111$ \\
Q-Learning        & $0{,}887$ & $0{,}061$ & $0{,}871$ & Extrap.\ de severidad ($0{,}762$) & $0{,}124$ \\
Q-Learning base   & $0{,}871$ & $0{,}074$ & $0{,}851$ & Extrap.\ de severidad ($0{,}627$) & $0{,}244$ \\
\bottomrule
\end{tabular}
\end{table}

\begin{figure}[H]
  \centering
  \includegraphics[width=\linewidth]{ind_01_desempeno_por_escenario.png}
  \caption[Desempe\~no medio por escenario (modelos
  individuales)]{Modelos individuales: desempe\~no medio $\bar r$ por
  modelo y escenario. La primera columna es la referencia.}
  \label{fig:heatmap-global}
\end{figure}

En la Figura~\ref{fig:extra-sev-skew}, con las secuencias ordenadas de
menor a mayor $\bar r$, el Transformer queda por encima o al nivel del
resto en los tres escenarios de severidad y en el de longitudes Poisson,
con la mayor ventaja en los peores casos de la distribuci\'on bimodal;
con secuencias largas lo supera el DQN recurrente, y con alta severidad y
longitud $10$ las curvas casi se superponen. Las
Figuras~\ref{fig:c-base}--\ref{fig:c-trf} muestran, para cada modelo, el
desempe\~no de cada secuencia de la referencia, su distribuci\'on y sus
estad\'isticos por bloque.

\begin{figure}[H]
  \centering
  \includegraphics[width=\linewidth]{ind_05_curvas_por_familia.png}
  \caption[Curvas por secuencia en seis escenarios (modelos
  individuales)]{Desempe\~no por secuencia, ordenado de menor a mayor,
  en seis escenarios de generalizaci\'on. Fila superior: severidad;
  fila inferior: longitud y conjunta. El trazo grueso corresponde al
  Transformer.}
  \label{fig:extra-sev-skew}
\end{figure}

\begin{figure}[H]
  \centering
  \includegraphics[width=0.64\linewidth]{PES_BASE_results.png}
  \caption[Resultados por secuencia: Q-Learning base]{Resultados por secuencia de Q-Learning base (\texttt{pes\_base}) en la referencia.}
  \label{fig:c-base}
\end{figure}

\vspace{-1.35\baselineskip} % entre figuras por modelo consecutivas
\begin{figure}[H]
  \centering
  \includegraphics[width=0.64\linewidth]{PES_QL_results.png}
  \caption[Resultados por secuencia: Q-Learning]{Resultados por secuencia de Q-Learning (\texttt{pes\_ql}) en la referencia.}
  \label{fig:c-ql}
\end{figure}

\vspace{-1.35\baselineskip} % entre figuras por modelo consecutivas
\begin{figure}[H]
  \centering
  \includegraphics[width=0.64\linewidth]{PES_DQL_results.png}
  \caption[Resultados por secuencia: Double Q-Learning]{Resultados por secuencia de Double Q-Learning (\texttt{pes\_dql}) en la referencia.}
  \label{fig:c-dql}
\end{figure}

\vspace{-1.35\baselineskip} % entre figuras por modelo consecutivas
\begin{figure}[H]
  \centering
  \includegraphics[width=0.64\linewidth]{PES_DQN_results.png}
  \caption[Resultados por secuencia: DQN]{Resultados por secuencia de DQN (\texttt{pes\_dqn}) en la referencia.}
  \label{fig:c-dqn}
\end{figure}

\vspace{-1.35\baselineskip} % entre figuras por modelo consecutivas
\begin{figure}[H]
  \centering
  \includegraphics[width=0.64\linewidth]{PES_RDQN_results.png}
  \caption[Resultados por secuencia: DQN recurrente]{Resultados por secuencia del DQN recurrente (\texttt{pes\_rdqn}) en la referencia.}
  \label{fig:c-rdqn}
\end{figure}

\vspace{-1.35\baselineskip} % entre figuras por modelo consecutivas
\begin{figure}[H]
  \centering
  \includegraphics[width=0.64\linewidth]{PES_A2C_results.png}
  \caption[Resultados por secuencia: A2C]{Resultados por secuencia de A2C (\texttt{pes\_a2c}) en la referencia.}
  \label{fig:c-a2c}
\end{figure}

\vspace{-1.35\baselineskip} % entre figuras por modelo consecutivas
\begin{figure}[H]
  \centering
  \includegraphics[width=0.64\linewidth]{PES_TRF_results.png}
  \caption[Resultados por secuencia: Transformer]{Resultados por secuencia del Transformer (\texttt{pes\_trf}) en la referencia.}
  \label{fig:c-trf}
\end{figure}

\subsection{Ensambles}\label{sec:res-ensembles}

La Tabla~\ref{tab:ensemble-stress} resume los seis ensambles y la
Tabla~\ref{tab:trf-vs-ens} compara cada uno con el Transformer individual;
como los dos grupos se evaluaron por separado, $d$ y $p$ se calculan con la
media, la desviaci\'on est\'andar y el n\'umero de secuencias de cada
escenario.

S\'olo el ensamble ponderado supera al Transformer: $0{,}939$
frente a $0{,}930$ en generalizaci\'on, mayor media en $18$ de los $21$
escenarios y un peor escenario $0{,}040$ m\'as alto. La diferencia es
significativa pero de tama\~no despreciable ($d = 0{,}19$, apenas por
debajo de un efecto peque\~no). En la referencia reduce en un $45\%$ la
distancia al \'optimo que deja Q-Learning,
$(\bar r - \bar r_{\mathrm{ql}})/(1 - \bar r_{\mathrm{ql}})$, frente al
$36\%$ del Transformer. Frente a los dem\'as ensambles, su ventaja es de tama\~no
medio sobre el consenso ($d = 0{,}58$) y despreciable sobre la compuerta
($d = 0{,}17$); ambas son significativas ($p < 10^{-5}$).

\begin{table}[H]
\centering
\caption[Desempe\~no de los ensambles]{Ensambles: media en la referencia,
media en generalizaci\'on, peor escenario y p\'erdida respecto de la
referencia.}
\label{tab:ensemble-stress}
\small
\setlength{\tabcolsep}{4pt}
\begin{tabular}{@{}lcclc@{}}
\toprule
\textbf{Ensamble} & \textbf{Ref.} & \textbf{Gen.} & \textbf{Peor escenario} & \textbf{P\'erdida} \\
\midrule
Ensamble ponderado        & $0{,}937$ & $0{,}939$ & Secuencias largas ($0{,}900$) & $0{,}037$ \\
Compuerta del Transformer & $0{,}928$ & $0{,}931$ & Secuencias largas ($0{,}861$) & $0{,}067$ \\
Consenso con prior        & $0{,}917$ & $0{,}922$ & Secuencias largas ($0{,}871$) & $0{,}046$ \\
Consenso                  & $0{,}889$ & $0{,}905$ & Secuencias largas ($0{,}826$) & $0{,}063$ \\
Voto suave                & $0{,}914$ & $0{,}902$ & Extrap.\ de severidad ($0{,}860$) & $0{,}054$ \\
Voto por acci\'on         & $0{,}914$ & $0{,}901$ & Extrap.\ de severidad ($0{,}859$) & $0{,}055$ \\
\bottomrule
\end{tabular}
\end{table}

\begin{table}[H]
\centering
\caption[Comparaci\'on de cada ensamble con el Transformer]{Cada ensamble
frente al Transformer individual en generalizaci\'on; los valores
positivos indican ventaja del ensamble.}
\label{tab:trf-vs-ens}
\small
\setlength{\tabcolsep}{4pt}
\begin{tabular}{@{}lccccc@{}}
\toprule
\textbf{Ensamble} & \textbf{Diferencia} & $d$ & \textbf{Efecto} & $p$ & \textbf{Mayor media en} \\
\midrule
Ensamble ponderado        & $+0{,}010$ & $0{,}19$  & despreciable & $< 10^{-6}$  & $18$ de $21$ \\
Compuerta del Transformer & $+0{,}001$ & $0{,}02$  & despreciable & $0{,}64$     & $15$ de $21$ \\
Consenso con prior        & $-0{,}008$ & $-0{,}15$ & despreciable & $< 10^{-3}$  & $4$ de $21$ \\
Consenso                  & $-0{,}025$ & $-0{,}39$ & peque\~no    & $< 10^{-23}$ & $3$ de $21$ \\
Voto suave                & $-0{,}027$ & $-0{,}41$ & peque\~no    & $< 10^{-24}$ & $2$ de $21$ \\
Voto por acci\'on         & $-0{,}028$ & $-0{,}42$ & peque\~no    & $< 10^{-26}$ & $2$ de $21$ \\
\bottomrule
\end{tabular}
\end{table}

La compuerta del Transformer obtiene un desempe\~no equivalente al del
Transformer ($d = 0{,}02$; $p = 0{,}64$): supera su media en $15$
escenarios, pero nunca por m\'as de $0{,}006$. Los otros cuatro ensambles
quedan por debajo del Transformer individual. El voto suave y el voto por
acci\'on, que s\'olo promedian o votan ponderando por la confianza, tienen
las medias m\'as bajas (efecto peque\~no, $d \approx -0{,}4$). Su peor
escenario es la extrapolaci\'on de severidad: el $41\%$ y el $42\%$ de sus
secuencias quedan por debajo de $0{,}8$
(Figura~\ref{fig:ensemble-extrapolation}) y su media queda por debajo de
la de tres de sus cuatro miembros: el Transformer ($0{,}996$), A2C
($0{,}929$) y DQN ($0{,}890$).

El mejor ensamble es el ponderado: tiene la mayor media en la
referencia ($0{,}937$) y en generalizaci\'on ($0{,}939$), el peor escenario
m\'as alto ($0{,}900$, secuencias largas) y la menor p\'erdida respecto de
la referencia ($0{,}037$) (Tabla~\ref{tab:ensemble-stress}). Es adem\'as el
ensamble con mayor media en $19$ de los $22$ escenarios; s\'olo lo superan
el consenso, en \texttt{len\_all\_short} y \texttt{joint\_low\_short}, y
la compuerta del Transformer, en \texttt{len\_poisson}
(Figura~\ref{fig:heatmap-ens}).

\begin{figure}[H]
  \centering
  \includegraphics[width=\linewidth]{ens_01_desempeno_por_escenario.png}
  \caption[Desempe\~no medio por escenario (ensambles)]{Ensambles:
  desempe\~no medio $\bar r$ por ensamble y escenario. La primera columna
  es la referencia.}
  \label{fig:heatmap-ens}
\end{figure}

\subsubsection{Efecto de las reglas fijas sobre la decisi\'on}\label{sec:res-freq}

Tres ensambles agregan reglas fijas al voto: el ensamble ponderado y el
consenso con \emph{prior} lo mezclan con un \emph{prior} de severidad y
aplican una cota de seguridad, y la compuerta del Transformer decide
cu\'ando ejecutar la acci\'on del Transformer; el voto suave, el voto por
acci\'on y el consenso no tienen reglas fijas. Para saber si la decisi\'on
de esos tres ensambles proviene del voto o de las reglas, se repiti\'o su
evaluaci\'on registrando, en cada decisi\'on, la
acci\'on antes y despu\'es de cada regla y la que prefiere el Transformer
(Ap\'endice~\ref{ap:orchestrator}). La repetici\'on reproduce las medias de
la evaluaci\'on principal con diferencias menores que el $0{,}01\%$ de la
escala de $\bar r$ (Tabla~\ref{tab:ens-freq}).

En el consenso con \emph{prior} decide el voto: el \emph{prior} cambia la
acci\'on votada en el $2\%$ de las decisiones de la referencia y en el
$4\%$ de las de generalizaci\'on, y el ensamble queda por debajo del
Transformer. En la compuerta decide el Transformer: su acci\'on se ejecuta
en el $97\%$ y el $95\%$ de las decisiones, lo que explica la equivalencia
de sus desempe\~nos. En el ensamble ponderado deciden el voto y el
\emph{prior}: con $\tau = 15$ las distribuciones de los miembros quedan
casi planas (en la referencia, la acci\'on m\'as probable de cada uno tiene
en promedio probabilidad $0{,}19$ y supera a la segunda por
$0{,}02$--$0{,}03$), y el \emph{prior} cambia la acci\'on votada en el
$41\%$ y el $33\%$ de las decisiones. La cota casi nunca act\'ua. Por lo
tanto, en los dos ensambles que igualan o superan al Transformer la
decisi\'on no proviene s\'olo del voto ponderado por confianza, sino de la
regla.

\begin{table}[H]
\centering
\caption[Frecuencia con que las reglas cambian la decisi\'on]{Porcentaje
de las decisiones con recursos disponibles en que ocurre cada evento.}
\label{tab:ens-freq}
\small
\setlength{\tabcolsep}{4pt}
\begin{tabular}{@{}llcc@{}}
\toprule
\textbf{Ensamble} & \textbf{Evento} & \textbf{Ref.} & \textbf{Gen.} \\
\midrule
Compuerta del Transformer & Sigue al Transformer ($g \ge 0{,}5$) & $96{,}9\%$ & $95{,}0\%$ \\
                          & Acci\'on final distinta de la del Transformer & $0{,}8\%$ & $2{,}7\%$ \\
\midrule
Consenso con prior & El \emph{prior} cambia la acci\'on votada & $2{,}4\%$ & $3{,}6\%$ \\
                   & La cota cambia la acci\'on & $0{,}0\%$ & $1{,}1\%$ \\
                   & Acci\'on final distinta de la del Transformer & $63{,}5\%$ & $57{,}9\%$ \\
\midrule
Ensamble ponderado & El \emph{prior} cambia la acci\'on votada & $40{,}9\%$ & $32{,}5\%$ \\
                   & La cota cambia la acci\'on & $0{,}0\%$ & $0{,}3\%$ \\
                   & Acci\'on final distinta de la del Transformer & $61{,}2\%$ & $48{,}1\%$ \\
\bottomrule
\end{tabular}
\end{table}

\subsection{Transformer frente al mejor ensamble}\label{sec:res-trf-vs-best}

El ensamble ponderado supera al Transformer en todas las medidas de la
Tabla~\ref{tab:trf-vs-best}: $+0{,}010$ en la media de generalizaci\'on,
con un efecto despreciable ($d = 0{,}19$; Tabla~\ref{tab:trf-vs-ens}),
$+0{,}040$ en el peor escenario y una p\'erdida m\'axima de poco m\'as de la
mitad. El Transformer tiene mayor media s\'olo en \texttt{len\_all\_short}
($0{,}936$ frente a $0{,}908$), en \texttt{joint\_low\_short} ($0{,}969$
frente a $0{,}942$) y, por menos de $0{,}001$, en \texttt{len\_poisson}.
En la extrapolaci\'on de severidad y en la conjunta, el ensamble ponderado
alcanza el \'optimo en todas las secuencias ($\sigma = 0$;
Figura~\ref{fig:ensemble-extrapolation}), pero todas sus decisiones
coinciden con las del Transformer que lo integra: la diferencia con el
Transformer individual proviene de que el ensamble recorta la severidad a
$9$ antes de pas\'arsela a sus miembros. La Figura~\ref{fig:c-ens} muestra
su desempe\~no por secuencia en la referencia, donde tiene la menor
dispersi\'on de todos los sistemas ($\sigma = 0{,}035$).

\begin{table}[H]
\centering
\caption[Transformer frente al ensamble ponderado]{Transformer individual
frente al ensamble ponderado, el mejor ensamble.}
\label{tab:trf-vs-best}
\small
\begin{tabular}{@{}lcc@{}}
\toprule
\textbf{Medida} & \textbf{Transformer} & \textbf{Ensamble ponderado} \\
\midrule
Media en la referencia ($\sigma$) & $0{,}927$ ($0{,}045$) & $0{,}937$ ($0{,}035$) \\
Media en generalizaci\'on         & $0{,}930$ & $0{,}939$ \\
Peor escenario (secuencias largas) & $0{,}860$ & $0{,}900$ \\
Mayor p\'erdida                    & $0{,}068$ & $0{,}037$ \\
Extrapolaci\'on de severidad       & $0{,}996$ & $1{,}000$ \\
Extrapolaci\'on conjunta           & $0{,}997$ & $1{,}000$ \\
Escenarios con mayor media (de $21$) & $3$ & $18$ \\
\bottomrule
\end{tabular}
\end{table}

\vspace{-\baselineskip} % entre flotantes consecutivos, sin el espacio extra
\begin{figure}[H]
  \centering
  \includegraphics[width=0.9\linewidth]{ens_06_curvas_extrapolacion.png}
  \caption[Curvas de los ensambles en los escenarios fuera de
  rango]{Desempe\~no por secuencia, ordenado de menor a mayor, de los seis
  ensambles en los escenarios fuera de rango; trazo grueso: ensamble
  ponderado; l\'ineas punteadas: media de cada ensamble.}
  \label{fig:ensemble-extrapolation}
\end{figure}

\vspace{-\baselineskip} % entre figuras consecutivas, sin el espacio extra
\begin{figure}[H]
  \centering
  \includegraphics[width=0.66\linewidth]{PES_ENS_results.png}
  \caption[Resultados por secuencia: ensamble ponderado]{Resultados por secuencia del ensamble ponderado (\texttt{pes\_ens}) en la referencia.}
  \label{fig:c-ens}
\end{figure}

\biblio % Needed for referencing to working when compiling individual subfiles - Do not remove
\end{document}
